\documentclass[12pt,a4paper]{article}
\usepackage[utf8]{inputenc}
\usepackage[polish]{babel}
\usepackage{amsmath}
\usepackage{amssymb}

\title{Algorytmy grafowe -- przeglad}
\author{Alesia Filinkova \and Diana Pelin}
\date{Maj 2026}

\begin{document}
\maketitle

\section{Wprowadzenie}

Graf to para uporzadkowana $G = (V, E)$, gdzie $V$ jest skonczonym
zbiorem wierzcholkow, a $E \subseteq V \times V$ zbiorem krawedzi.
Algorytmy grafowe stanowia jedna z podstawowych galezi informatyki
i znajduja zastosowanie wszedzie, gdzie modelujemy relacje miedzy
obiektami: w sieciach komputerowych, planowaniu trasy, kompilatorach
oraz w analizie sieci spolecznych.

\section{Przeszukiwanie wszerz}

Algorytm BFS odwiedza wierzcholki w kolejnosci ich odleglosci od
wierzcholka startowego. Wykorzystuje strukture kolejki FIFO i ma
zlozonosc $O(V + E)$. Dla grafu nieskierowanego BFS pozwala wyznaczyc
najkrotsze sciezki w sensie liczby krawedzi.

\section{Algorytm Dijkstry}

Dijkstra znajduje najkrotsze sciezki ze zrodla do wszystkich
pozostalych wierzcholkow w grafie z nieujemnymi wagami krawedzi.
Implementacja oparta o kopiec binarny ma zlozonosc $O((V + E) \log V)$.
Algorytm zaklada, ze waga $w(u, v) \geq 0$ dla kazdej krawedzi.

\section{Drzewo rozpinajace}

Algorytmy Prima i Kruskala wyznaczaja minimalne drzewo rozpinajace
w grafie spojnym. Kruskal wykorzystuje strukture zbiorow rozlacznych,
a Prim -- kopiec priorytetowy. Oba osiagaja zlozonosc $O(E \log V)$.

\end{document}
